\documentclass[12pt]{article}
\usepackage{mathtools}
\usepackage{booktabs}
\usepackage{enumitem}
\usepackage{microtype}
\usepackage{setspace}
\usepackage{hyperref}
\usepackage{natbib}
\usepackage{amsthm}
\usepackage{lastpage}

\newcommand{\bb}{\mathbf}
\newcommand{\eps}{\varepsilon}
\newcommand{\KL}{\operatorname{KL}}
\newcommand{\E}{\mathbb{E}}
\newcommand{\R}{\mathbb{R}}
\newcommand{\Prob}{\mathbb{P}}
\newcommand{\Var}{\operatorname{Var}}
\newcommand{\Iden}{\operatorname{I}}
\newcommand{\PiW}{\Pi_{W}}
\newcommand{\argmin}{\operatorname*{arg\,min}}

\newtheorem{theorem}{Theorem}[section]
\newtheorem{lemma}[theorem]{Lemma}
\newtheorem{proposition}[theorem]{Proposition}
\newtheorem{corollary}[theorem]{Corollary}
\newtheorem{conjecture}[theorem]{Conjecture}
\newtheorem{definition}[theorem]{Definition}
\newtheorem{assumption}[theorem]{Assumption}
\newtheorem{prediction}[theorem]{Prediction}
\theoremstyle{remark}
\newtheorem{remark}{Remark}

\begin{document}
\begin{center}{\LARGE\bf Supporting Information}\end{center}
\begin{center}{\large\bf A Rate--Distortion Framework for Language-Model Scaling Laws}\end{center}
\doublespacing

This Supporting Information contains the full proofs, detailed numerical
diagnostics, and bootstrap tables referenced in the main manuscript.  All
labeling follows the main text; cross-references use the same keys.


\section{Detailed toy-model numerical tables}
\label{app:numerics_tables}

The following tables provide the detailed numerical diagnostics summarized in
Section~10 of the main manuscript.

\paragraph{Two-token Case I/II.}
\label{numerics:twotoken}
\begin{center}
\begin{tabular}{ccc}
\toprule
Channel & $\E[\text{MSE}]$ & $\E[\text{KL}]$ \\
\midrule
bounded $f=\frac12+0.4\,\operatorname{tri}$ & $W^{-3.00}$ & $W^{-2.99}$ \\
degenerate $f=\operatorname{tri}$ & $W^{-3.00}$ & $W^{-2.01}$ \\
\bottomrule
\end{tabular}
\end{center}
Case~I (bounded) preserves the exponent; Case~II (boundary touch) shifts the
KL exponent by $+1$ in $W$, i.e., $\alpha_{\rm KL}=\alpha_{L^2}-\gamma$.

\paragraph{V-token generalization.}
\label{numerics:vtoken}
$V=3$, $P_1=(1-\rho)\operatorname{tri}$, $P_2=(1-\rho)(1-\operatorname{tri})$,
$P_3=\rho$, renormalized projection:
\begin{center}
\begin{tabular}{cc}
\toprule
$\rho$ & exponent (W = 32..256) \\
\midrule
0.10 & $-2.01\ -2.00\ -1.99\ -1.98$ \\
0.30 & $-2.01\ -2.00\ -1.99\ -1.98$ \\
0.50 & $-2.01\ -2.00\ -1.99\ -1.98$ \\
\bottomrule
\end{tabular}
\end{center}
Exponent $-2.00$ exactly, $\rho$-independent: the Case~II boundary shift
survives $V=3$, no cancellation among the $V-1$ positive weight terms.

\paragraph{Parameterization.}
Model~B (logit-space truncation, channel touches zero): the logit
$\log(1-\operatorname{tri})$ has a log-singularity at the triangle apex; its
Fourier coefficients verify the $1/k$ tail (Appendix~\ref{app:logit}):
\begin{center}
\begin{tabular}{cccc}
\toprule
$k$ & 1 & 3 & 7 \\
\midrule
$|c_k|$ & 1.20 & 0.38 & 0.17 \\
$k\,|c_k|$ & 1.20 & 1.14 & 1.19 \\
\bottomrule
\end{tabular}
\end{center}
($k|c_k|\approx1.13$--$1.33$ throughout), and $E[\KL]$ flattens toward $-1$
($-1.39\ -1.30\ -0.88\ -0.55\ -0.30$ at $W=32,64,128,256,512$) before the
numerical floor. Model~C (logit space, strictly positive channel): clean
Case~I ($-3.08\ -3.07\ -2.81$ until the quadrature floor at $W\approx128$).
The two models bracket the real situation: real near-degenerate logits sit
between them (Remark~\ref{rem:logit}).

\paragraph{Additivity.}
Squared-loss cross term $=\,+1.05\times10^{-9}$ across every $(W,D)$ row
(identical MC residual across rows---exactly zero in exact arithmetic).
KL cross/total: bounded channel $-5.0\times10^{-3}\to-8.7\times10^{-5}$ for
$W=8\to64$; degenerate channel $-0.12\to-7.1\times10^{-3}$. Interior KL cross
in data expectation, $D=10^5$, $M=30$ draws (Theorem~\ref{thm:additive}):
\begin{center}
\begin{tabular}{ccc}
\toprule
$W$ & $\E[\delta_b\delta_v]$ & $|r|/\eps$ \\
\midrule
8  & $+1.1\times10^{-6}$ & $8.4\times10^{-3}$ \\
16 & $+1.3\times10^{-7}$ & $6.3\times10^{-3}$ \\
32 & $-7.1\times10^{-8}$ & $1.6\times10^{-2}$ \\
64 & $-1.5\times10^{-10}$ & $2.3\times10^{-2}$ \\
\bottomrule
\end{tabular}
\end{center}
The quadratic cross vanishes in expectation (the rows are MC noise), and the
remainder is the $O(\eps^{3/2})$ cubic correction, $o(\eps)$. The KL/MSE ratio
is $\approx8$ for the bounded channel, constant in $W$, consistent with the
weighted-$L^2$ bridge (Lemma~\ref{lem:bridge}).

\paragraph{Saturation.}
Wiener-filter risk of Theorem~\ref{thm:saturation}, bounded channel,
$2s/\beta_{\rm reg}=4$ (predicted D-exponent $-3/4$):
\begin{center}
\begin{tabular}{ccc}
\toprule
fixed $D$, $W=16\to512$ & risk change & $<1\%$ \\
fixed $W=512$ & D-exponent & $-0.83,-0.80,-0.78,-0.77\to-3/4$ \\
oracle Wiener, $D=2\times10^5$ & KL & $4.7/5.2/6.4/5.6\times10^{-5}$ ($W=16..256$)\\
plain ERM, same & KL & $1.7\times10^{-4}\to1.3\times10^{-3}$ (grows) \\
degenerate channel & KL & $2.1/2.4/1.9/1.7\times10^{-4}$ ($W=16..256$)\\
\bottomrule
\end{tabular}
\end{center}
The variance saturates in $W$ for the self-truncating estimator, including for
the degenerate (boundary-touching) channel; unregularized ERM grows.
Figure~\ref{fig:synth} shows the corresponding spectral diagnostics.


\clearpage
\appendix
\section{Full proofs}
\label{app:proofs}

\subsection{The exact KL expansion}
\label{app:klseries}
For binary channels, let $g-f=\delta$ pointwise. Then
\begin{align}
  \KL(f\,\|\,f+\delta)
  &= -f\log\!\Bigl(1+\frac{\delta}{f}\Bigr)
   -(1-f)\log\!\Bigl(1-\frac{\delta}{1-f}\Bigr) \notag\\
  &= \frac{\delta^2}{2}\Bigl(\frac1f+\frac1{1-f}\Bigr)
   + \frac{\delta^3}{3}\Bigl(\frac1{(1-f)^2}-\frac1{f^2}\Bigr)
   + \frac{\delta^4}{4}\Bigl(\frac1{f^3}+\frac1{(1-f)^3}\Bigr)+\cdots
  \label{eq:klseries}
\end{align}
Term by term: $-f\log(1+\delta/f)=-\delta+\delta^2/(2f)-\delta^3/(3f^2)+\cdots$
and $-(1-f)\log(1-\delta/(1-f))=\delta+\delta^2/(2(1-f))+\delta^3/(3(1-f)^2)+\cdots$;
the first-order terms cancel (KL is first-order flat in $\delta$), the
second-order term is $\delta^2/(2f(1-f))$ (the conditional inverse-Bernoulli
variance of Remark~\ref{rem:conditional}), and the remainder is
$O(\delta^3)$ with coefficients bounded in the interior. This exact series
sits behind Lemma~\ref{lem:bridge} and is used again in the proof of
Theorem~\ref{thm:additive}.

\subsection{Proof of Theorem~\ref{thm:modelN} (two-sided)}
\label{app:modelN}
\textit{Upper bound.} In the small-error regime,
$\eps_N\asymp\frac12\E_x\sum_y(\PiW f-f)_y^2/f_y$ by
Lemma~\ref{lem:bridge}. The interior bound $f_y\ge\delta$ gives
$\eps_N\le\frac1{2\delta}\E_x\sum_y(\PiW f-f)_y^2$. By Parseval and the
$s$-source condition $\langle f_y,e_i\rangle^2\le i^{-2s/\beta_{\rm reg}}$,
\begin{equation*}
  \E_x\sum_y(\PiW f-f)_y^2
  = \sum_y\sum_{i>W}\langle f_y,e_i\rangle^2
  \le \sum_y\sum_{i>W}i^{-2s/\beta_{\rm reg}}
  \lesssim W^{-(2s/\beta_{\rm reg}-1)},
\end{equation*}
where the last bound uses $\sum_{i>W}i^{-\mu}\asymp W^{-(\mu-1)}$ for
$\mu>1$ (Euler--Maclaurin). \textit{Lower bound.} The operator
$T=\sum_y K_X^{2s}$ has eigenvalues $i^{-2s/\beta_{\rm reg}}$, so
$\sum_{i>W}\langle f_y,e_i\rangle^2=\langle f_y,T_{>W}f_y\rangle$ has exactly
the same tail for the mode-optimal token family (achieved at any $y$ with
$\|v_y\|>0$); the interior weight is bounded below by $1/(1-\delta)$-type
terms, giving the two-sided $\asymp$.

\subsection{Proof of Proposition~\ref{prop:boundary}}
\label{app:boundary}
With $f\sim au$ near the boundary point and truncation error $c\ne0$ (the
projection cannot reproduce the cusp; generically $|c|\sim
W^{-(s/\beta_{\rm reg}-1)}$), the KL series \eqref{eq:klseries} gives
$\KL\approx\frac{c^2}{2}\bigl(\frac1f+\frac1{1-f}\bigr)=\frac{c^2}{2au}+O(1)$
for $1/W\ll u\ll1$, and $\KL=O(1/W)$ for $u\lesssim 1/W$ (the projection
resolves the cusp locally below wavelength $1/W$). The mean excess integrates
the two regimes against $u\sim$ uniform:
\begin{equation*}
  \eps_N^{\rm bdy}
  \sim \int_{0}^{1}\KL(u)\,du
  \sim \int_{1/W}^{1}\frac{c^2}{4u}\,du + O(1/W^2)
  = \frac{c^2}{4}\log W + O(1/W^2),
\end{equation*}
which is \eqref{eq:bdyint}. For the triangle channel the first odd-mode
coefficient is $c=-4/(\pi^2 W)$ (closed form), so
$\eps_N^{\rm bdy}\asymp W^{-2}$ up to an unresolved $\log W$ factor. The
general boundary-approach rate $u^p$ is stated as an open problem
(Appendix~\ref{app:open}).

\subsection{Proof of Proposition~\ref{prop:vtoken}}
\label{app:vtoken}
For $V\ge2$ tokens, in the small-error regime
$\KL\approx\frac12\sum_y\delta_y^2/f_y$. Near a boundary point where token
$y$ has $f_y\sim a_y u$ and truncation error $\delta_y=c_y+O(u)$:
\begin{equation*}
  \frac{\delta_y^2}{2f_y}\sim\frac{c_y^2}{2a_y u}\quad(u\to0),
\end{equation*}
with all terms positive. Integrating gives a boundary contribution
$\bigl(\sum_y c_y^2/(2a_y)\bigr)\log W$: the terms add, dominated by the
tokens with $c_y\ne0$, and there is no cancellation. Tokens bounded away from
zero contribute only to the interior (Case~I) part.

\subsection{Logit-space spectral tail (Proposition~\ref{prop:logit})}
\label{app:logit}
If the channel touches zero, the logit $\ell(x)=\log f(x)$ has a
log-singularity, $\ell(x)=\ell_0(x)+a\log|2\sin(\pi(x-x_0))|$ locally. The
classical series $\log|2\sin(\pi x)|=-\sum_{k\ge1}\frac{\cos(2\pi kx)}{k}$
shows the Fourier coefficients of a log-singularity decay as $1/k$, so the
truncated-logit squared tail is
$\sum_{k>W}|c_k|^2\asymp\sum_{k>W}k^{-2}\asymp W^{-1}$. After softmax the
channel perturbation is $\delta_p=p\,(\delta_\ell-\E_p\delta_\ell)$ (first
order in the logit error), so
$\E_x[\KL]\asymp\E_x[\delta_p^2/p]\asymp\E_x[p\,\delta_\ell^2]\asymp
W^{-1}$ in the interior, versus $W^{-2}$ in probability space: the slowdown
is amplified by the logit parameterization. For strictly positive channels
the logit is smooth and its spectrum lies in the same class as the channel,
so Case~I survives (verified: $W^{-3.08},W^{-3.07},W^{-2.81}$).

\subsection{Exact integrals for the toy models}
\label{app:exact}
The triangle channel $f(x)=2|x-\tfrac12|$ has the exact Fourier series
$\operatorname{tri}=\frac12-\frac4{\pi^2}\sum_{k\ \text{odd}}
\frac{\cos(2\pi kx)}{k^2}$ (integrating by parts twice; the only singular
object is the derivative jump at $x=\tfrac12$). Hence the $W$-truncation error
and its mean square are exactly
\begin{equation*}
  g_W(x)=-\frac4{\pi^2}\sum_{\substack{k>W\\ k\ \text{odd}}}
  \frac{\cos(2\pi kx)}{k^2},
  \qquad
  \E[g_W^2]=\frac{16}{\pi^4}\sum_{\substack{k>W\\ k\ \text{odd}}}k^{-4}
  =\frac{8}{3\pi^4}W^{-3}\bigl(1+o(1)\bigr),
\end{equation*}
where the closed form uses $\sum_{\text{odd}\,k>W}k^{-4}=\frac16W^{-3}(1+o(1))$
(numerically the ratio reaches $1.000$ at $W=256$). Two further exact
evaluations anchor the boundary mechanism. First, the truncation error at the
zero-touch point,
\begin{equation*}
  g_W(\tfrac12)=\frac4{\pi^2}\sum_{\substack{k>W\\ k\ \text{odd}}}k^{-2}
  =\frac{2}{\pi^2}W^{-1}\bigl(1+o(1)\bigr)\quad
  (\text{ratio }0.999\text{ at }W=128),
\end{equation*}
is $O(W^{-1})$, i.e.\ the residual is polynomially larger where the channel
touches zero than in the bulk (there it is $O(W^{-2})$). Second, the KL is
computed exactly (not through the quadratic bridge): numerically
$\E[\KL]\cdot W^2=0.0265$ stable to three decimals across $W=16$--$1024$,
reproducing the exponent $-2.00$ and the interior/boundary split of
Proposition~\ref{prop:boundary} (at $W=256$ the boundary layer, of measure
$4\cdot10^{-2}$, contributes $1.0\cdot10^{-5}$ against $6.7\cdot10^{-10}$ from
the interior).

The exact integration exposes a microscopic mechanism that the quadratic
bridge \eqref{eq:bridge} conceals. Near the apex $u=x-\tfrac12$, $f\sim2|u|$ and
the quadratic weight $g^2/(2f(1-f))\sim g^2/(4|u|)$ diverges, so a naive
second-order evaluation overcounts the apex by an order of magnitude
(numerically $\approx10\times$). The exact two-token KL, however, saturates in
the layer $|u|\lesssim g$ where the perturbation exceeds the channel value:
there the leading term is $g_W^2/2=O(W^{-2})$ rather than the divergent
$g_W^2/(2f)$. The finite $W^{-2}$ exponent is produced by integrating this
saturated profile; the crossover scale $|u|\sim W^{-1}$ is exactly the
self-truncation scale of Theorem~\ref{thm:saturation}. In the interior the
two agree, which is why Case~I (\S\ref{sec:modelsize}, bounded channels) shows no
such correction.

\emph{No closed-form coefficient for bounded channels.} For the bounded model
$f=\tfrac12+0.4\operatorname{tri}$, the bridge predicts
$\E[\KL]/\E[g^2]\to\frac12\E[\frac1{f(1-f)}]=\frac52\ln3\approx2.75$, but the
measured ratio rises $3.42\to3.76$ as $W$ grows rather than converging. The
reason is that the weight $1/(f(1-f))$ breaks basis orthogonality:
$\E[g^2/f(1-f)]$ carries cross terms $\sum_{k\ne l}c_kc_l\int\varphi_k
\varphi_l/f(1-f)$, which are not subleading at accessible $W$. The theory
therefore predicts exponents, not prefactors; this is why all quantitative
statements in \S\ref{sec:numerics} are log-ratios.

\subsection{Proof of Theorem~\ref{thm:saturation}}
\label{app:minimax}
\textit{Achievability.} Over the class $M_W$ consider shrinkage estimators
$\hat f_i=w_i\hat c_i$. The risk is
\begin{equation*}
  R(w)=\sum_{i\le W}\E(w_i\hat c_i-c_i)^2
  =\sum_{i\le W}\Bigl[w_i^2\frac{\sigma_i^2}{D}+(1-w_i)^2 c_i^2\Bigr],
\end{equation*}
minimized termwise at $w_i=c_i^2/(c_i^2+\sigma_i^2/D)$, giving the oracle risk
\begin{equation*}
  R^*=\sum_{i\le W}\frac{c_i^2\,(\sigma_i^2/D)}{c_i^2+\sigma_i^2/D},
\end{equation*}
which is \eqref{eq:minimax}. \textit{Threshold.} With
$|c_i|\sim i^{-s/\beta_{\rm reg}}$ and $\sigma_i^2\approx\sigma^2$, the
condition $c_i^2=\sigma^2/D$ gives $i=W^*(D)\sim D^{\beta_{\rm reg}/(2s)}$;
for $W\ge W^*$, modes $i>W^*$ carry weight $w_i\lesssim1$ and their
contribution is at most their bias $c_i^2$, summable to
$\sum_{i>W^*}c_i^2\sim W^{*-(2s/\beta_{\rm reg}-1)}$, which is a constant
fraction of the leading term; hence the risk is $W$-independent and equals
\eqref{eq:varstar}. \textit{Lower bound.} The sequence-model minimax bound
(the Pinsker/Efroimovich result, cf.\ \citet{tsybakov2009}, ch.~3, and
\citet{pinsker1980}) shows no estimator over the $\ell^2$ body beats the
oracle risk up to
$(1+o(1))$; a hard threshold at level $\sigma\sqrt{2\log W}/\sqrt D$ achieves
the same rate with logarithmic losses absorbed in the $(1+o(1))$. The
transfer to the KL excess uses Lemma~\ref{lem:bridge}: in the interior the
weight is bounded and mode-independent to leading order, so the exponent is
unchanged.

\subsection{The telescope (Lemma~\ref{lem:telescoping})}
\label{app:telescoping}
With $L_n=H_n+\eps_n$ and $L_{n-1}=H_{n-1}+\eps_{n-1}$, subtraction gives
\eqref{eq:tele} immediately. Summing over $n=1,\dots,n^*(P)$:
\begin{equation*}
  L_{n^*}(P)-L_0(P)
  =\bigl[H_{n^*}-H_0\bigr]+\bigl[\eps_{n^*}(P)-\eps_0(P)\bigr],
\end{equation*}
and normalizing by the horizon count and passing to the AR average yields
\eqref{eq:telefinal} with the boundary term $H_{n^*}-H_\infty$ and the
estimation term $\eps_{n^*}(P)-\eps_0(P)$. Substituting
$H_{n^*}-H_\infty\sim n^{*\,-\gamma_{\rm ent}}$ and
$n^*\sim P^{1/(2\beta_{\rm corr})}$ gives the data exponent
\eqref{eq:alphaD} whenever the estimation term decays faster (the dominance
condition, Remark~\ref{rem:dominance}).

\subsection{Compute-optimality calculus}
\label{app:compute}
Minimize $\eps(N,D)$ subject to $C=6ND$.

\emph{Source-limited regime,} $\eps=A N^{-\alpha_N}+B D^{-\alpha_D}$. With
$D=C/(6N)$, the Lagrangian critical point satisfies
\begin{align*}
  \frac{\partial}{\partial N}:&&
  -\alpha_N A N^{-\alpha_N-1} &+ \lambda\,6D =0,\\
  \frac{\partial}{\partial D}:&&
  -\alpha_D B D^{-\alpha_D-1} &+ \lambda\,6N =0,
\end{align*}
so $\frac{D}{N}=\frac{\alpha_N A}{\alpha_D B}N^{-\alpha_N-1}D^{\alpha_D+1}$,
i.e., $N^{\alpha_N}\propto D^{\alpha_D}$, giving
$N^*\propto D^{\alpha_D/\alpha_N}$ and, from $ND=C$, the two forms in
\eqref{eq:sourceopt}.

\emph{Resolution-limited regime,} $\eps=A N^{-\alpha_N}+C_0 N^{\gamma}/D$
(with $C_0$ the variance prefactor). Substituting $D=C/(6N)$,
$\eps(N)=A N^{-\alpha_N}+\frac{6C_0}{C}N^{\gamma+1}$, minimized at
$\alpha_N A N^{-\alpha_N-1}=\frac{6C_0}{C}(\gamma+1)N^{\gamma}$, giving
\begin{equation*}
  N^*\propto C^{1/(\alpha_N+\gamma+1)},\qquad
  D^*\propto C^{(\alpha_N+\gamma)/(\alpha_N+\gamma+1)},
\end{equation*}
and eliminating $C$: $N^*\propto D^{1/(\alpha_N+\gamma)}$, the superlinear
slope of \eqref{eq:resolopt}.

\subsection{The projection--realization gap}
\label{app:gap}
Let $\Pi_W p$ be the orthogonal projection of the true channel onto the
resolved subspace (the object Theorem~\ref{thm:modelN} analyzes) and $p_W$ the
channel a real model actually realizes. Write the realized error as
$p-p_W=(p-\Pi_W p)+(\Pi_W p-p_W)$; the first term is the
\emph{projection} (bias) error and the second the \emph{realization} (model
capacity plus optimizer) error. Applying the exact KL series
(Appendix~\ref{app:klseries}) and $2ab\le a^2+b^2$ gives, to leading order,
\begin{equation*}
  \E[\KL(p\|p_W)]
  \;\lesssim\;\E[\KL(p\|\Pi_W p)]\;+\;\tfrac12\,\E_w\bigl[\,\|\Pi_W
  p-p_W\|^2\bigr],
\end{equation*}
where $\E_w$ is the weighted $L^2$ expectation with weight
$1/(p(1-p))$ and the cross term $\int (p-\Pi_Wp)(\Pi_Wp-p_W)\,w$ vanishes in
expectation when the realization is uncorrelated with the projection residual
at the resolved scale. If the realization error decays with effective
resolution at rate $\alpha_{\rm opt}$, then
$\alpha_{\rm real}\ge\min(\alpha_N,\alpha_{\rm opt})$: the projection
analysis is an upper bound on the realized error, tight whenever
$\alpha_{\rm opt}\ge\alpha_N$. Assumption~\ref{ass:fastlearn} asserts exactly
this (the optimizer tracks the resolved subspace), so under A5
$\alpha_{\rm real}=\alpha_N$. Conversely, when the channel touches zero the
\emph{parameterization} gap dominates: a probability-space projection
(Eq.~\ref{eq:bridge}) and a logit-space truncation are different objects, and
the latter is slower ($W^{-1}$ vs.\ $W^{-2}$, Proposition~\ref{prop:logit}).
The realized exponent can then be strictly worse than the projection bound
even with $\alpha_{\rm opt}$ large, because the logit-space realization is not
the probability-space projection. This is the precise sense in which
Theorem~\ref{thm:modelN} is an upper bound rather than an identity
(Section~\ref{sec:discussion}).

\subsection{Open problems}
\label{app:open}
(i) The general boundary-approach rate $u^p$: for channels reaching the
boundary as $u^p$, the boundary exponent changes (the triangle is $p=1$; the
general formula is open). (ii) The relation between $\beta_{\rm corr}$ and
$\beta_{\rm reg}$ through the Toeplitz spectral density of a stationary
process, which would collapse the two exponents to one corpus-statistics set.
(iii) Whether the $W^{-2}$ boundary term carries $\log W$: not resolved at
accessible $W$, cosmetic. (iv) A defense of $\gamma=\frac12$ for transformers
at all scales (the vocabulary embedding term $2\mathcal{V}d$ lowers the
effective $\gamma$ at small $N$).

\section{Numerical setup}
\label{app:numerics}
All channels are built from $\operatorname{tri}(x)=2|x-\tfrac12|$ with exact
Fourier series $\operatorname{tri}=\frac12-\frac4{\pi^2}\sum_{k\ \text{odd}}
\frac{\cos(2\pi kx)}{k^2}$, so truncation tails are closed form and free of
quadrature noise. Projection estimators: $\hat c_i=\frac1D\sum_j
y_j e_i(x_j)$ with $y_j\sim\text{Bernoulli}(f(x_j))$ (D draws), renormalized
per Proposition~\ref{prop:vtoken}. KL integrated exactly. Exponents are
log-ratios across $W$; $M=30$--$100$ independent draws averaged where data
expectations appear. Configurations: Case~I/II at $W=8$--$512$, exact
coefficients; V-token at $W=32$--$256$, $D=10^5$; saturation at $D=10^4$--
$2\times10^5$, oracle Wiener weights $w_i=c_i^2/(c_i^2+\sigma_i^2/D)$; the
additivity runs fix the draw sequence across $W$ so that the cross term is
exactly comparable across rows (the identical $10^{-9}$ residual is the
roundoff of a common deterministic quantity). All scripts and data
generators are released with the paper in the project repository
(\url{https://github.com/Shahul9570/Research.git}).

\section{Bootstrap confidence intervals on spectral quantities}
\label{app:bootstrap}
We assess the sampling stability of the eigenvalue-decay slopes and
$S_i^2$ slopes by re-running the exact eigendecomposition
(Section~\ref{sec:empirical}) with $10$ independent random window subsamples
($10^6$ windows each) for each of two corpora (prose, code) and two window
sizes ($W=4$, $W=8$). The same top-$1000$ token restriction and fitting
procedure are used throughout. Percentile $95\%$ confidence intervals are
reported.

\paragraph{Eigenvalue slopes.}
\begin{center}
\small
\begin{tabular}{lcccc}
\toprule
& \multicolumn{2}{c}{$[1,10]$} & \multicolumn{2}{c}{$[10,50]$} \\
\cline{2-3} \cline{4-5}
Config & slope & $\beta_{\rm reg}$ & slope & $\beta_{\rm reg}$ \\
\midrule
prose $W\!=\!4$ & $-0.194$ $[-0.197,-0.192]$ & $5.15$ & $-0.765$ $[-0.767,-0.763]$ & $1.31$ \\
prose $W\!=\!8$ & $-0.129$ $[-0.131,-0.127]$ & $7.81$ & $-0.631$ $[-0.633,-0.629]$ & $1.58$ \\
code $W\!=\!4$ & $-0.443$ $[-0.444,-0.441]$ & $2.26$ & $-0.724$ $[-0.727,-0.722]$ & $1.38$ \\
code $W\!=\!8$ & $-0.195$ $[-0.196,-0.193]$ & $5.13$ & $-0.700$ $[-0.702,-0.698]$ & $1.43$ \\
\midrule
& \multicolumn{2}{c}{$[50,200]$} & \multicolumn{2}{c}{$[200,800]$} \\
\cline{2-3} \cline{4-5}
Config & slope & $\beta_{\rm reg}$ & slope & $\beta_{\rm reg}$ \\
\midrule
prose $W\!=\!4$ & $-1.386$ $[-1.392,-1.382]$ & $0.72$ & $-0.843$ $[-0.846,-0.836]$ & $1.19$ \\
prose $W\!=\!8$ & $-1.176$ $[-1.178,-1.173]$ & $0.85$ & $-1.150$ $[-1.152,-1.148]$ & $0.87$ \\
code $W\!=\!4$ & $-0.957$ $[-0.959,-0.954]$ & $1.04$ & $-0.988$ $[-0.994,-0.985]$ & $1.01$ \\
code $W\!=\!8$ & $-0.815$ $[-0.817,-0.813]$ & $1.23$ & $-1.041$ $[-1.044,-1.039]$ & $0.96$ \\
\bottomrule
\end{tabular}
\end{center}
The implied $\beta_{\rm reg}$ ranges from $0.72$ to $7.81$ across
configurations and fit windows. The CIs are tight (bootstrap std $<0.004$ for
all eigenvalue slopes), confirming the range-dependence is not sampling noise.

\paragraph{Channel-smoothness slopes.}
\begin{center}
\small
\begin{tabular}{lccc}
\toprule
Config & $[3,10]$ & $[10,40]$ & $[40,70]$ \\
\midrule
prose $W\!=\!4$ & $-5.90$ $[-6.95,-4.16]$ & $-0.06$ $[-0.62,+0.39]$ & $+0.01$ $[-1.06,+1.21]$ \\
prose $W\!=\!8$ & $-9.10$ $[-11.70,-7.63]$ & $+3.28$ $[+2.96,+3.96]$ & $-0.15$ $[-0.89,+0.97]$ \\
code $W\!=\!4$ & $-8.32$ $[-10.01,-6.26]$ & $+1.82$ $[+1.06,+2.30]$ & $+0.32$ $[-2.18,+0.53]$ \\
code $W\!=\!8$ & $-8.04$ $[-10.34,-5.08]$ & $+2.46$ $[+1.96,+2.84]$ & $+1.07$ $[+0.03,+2.36]$ \\
\bottomrule
\end{tabular}
\end{center}
$S_i^2$ slopes are non-monotonic across all configurations: steep negative at
$[3,10]$, flat or positive at $[10,40]$, and variable at $[40,70]$. The CIs
are wider (bootstrap std $0.3$--$1.7$) but the non-monotonic pattern is
consistent across seeds.

\paragraph{Entropy and correlation statistics.}
From the full $50$M-token corpus: $H_1^{\rm MM}$--$H_5^{\rm MM}=5.994$,
$4.421$, $3.360$, $2.308$, $1.428$ bits/token; $\gamma_{\rm ent}=0.05$
(degenerate, fit returns $E\to-\infty$); $\beta_{\rm corr}=0.42$.

% ------------------------------------------------------------------

\end{thebibliography}

\end{document}
